% Level 2, Week 12 homework -- reflection and holiday menu
\homeworkweek{12}{Showcase and Celebration}{15--20 minutes for the reflection. The menu is for the holidays: pick whatever you like.}
\namedateline

\begin{core}[My Level 2 reflection (about 10 minutes)]
\begin{enumerate}
  \item The part of my project I am proudest of is… \answerlines{1}
  \item The hardest bug I met in Level 2 was… and I fixed it by… \answerlines{2}
  \item One piece of feedback from the gallery walk that I agree with: \answerlines{1}
  \item Something I learned about working with a partner: \answerlines{1}
  \item In Level 3 I am most curious about (git save points, classes, Maze Muncher, something else?): \answerlines{1}
\end{enumerate}
\end{core}

\begin{challenge}[Version 2 plan]
Plan the next version of your project. Write three new features. For each one: the function you
would write, what goes in, what comes out, and one test for it.
\answerspace{4cm}
\end{challenge}

\begin{creative}[Holiday \enquote{keep coding} menu]
Pick one, two, or none: holidays are for resting too. Everything works offline, with no accounts.
\begin{itemize}
  \item \textbf{Level up:} add one wish from your feedback cards to your project.
  \item \textbf{Art jam:} make a generative art picture with a palette list, \texttt{random} and a function. If you met \texttt{random.seed} in the Art Jam, write down your favorite seed.
  \item \textbf{Secret messages:} use your Caesar cipher program to send a coded note to someone at home.
  \item \textbf{Teach it:} show a family member how to play your game, then how one function works.
  \item \textbf{Remix a classic:} turn a past exercise (pixel art, Etch-a-Sketch, Word Guess) into something new. Add a credit comment.
  \item \textbf{Unplugged:} design a board game whose rules are a game loop: what happens every turn?
\end{itemize}
What I chose and what I made: \answerlines{2}
\end{creative}

\begin{selfcheck}
  \item I answered all five reflection questions.
  \item I kept my certificate and my workbook somewhere safe.
  \item If I shared my project, it was with family and showed only my nickname.
  \item I took breaks from the screen (stretch, and look at something about 20 feet (6 meters) away for 20 seconds).
\end{selfcheck}
